%%%%%%%%%%%%%%%%%%%%%%%%%%%%%%%%%%%%%%%%%
% Journal Article
% LaTeX Template
% Version 2.0 (February 7, 2023)
%
% This template originates from:
% https://www.LaTeXTemplates.com
%
% Author:
% Vel (vel@latextemplates.com)
%
% License:
% CC BY-NC-SA 4.0 (https://creativecommons.org/licenses/by-nc-sa/4.0/)
%
% NOTE: The bibliography needs to be compiled using the biber engine.
%
%%%%%%%%%%%%%%%%%%%%%%%%%%%%%%%%%%%%%%%%%

%----------------------------------------------------------------------------------------
%	PACKAGES AND OTHER DOCUMENT CONFIGURATIONS
%----------------------------------------------------------------------------------------

\documentclass[
	a4paper, % Paper size, use either a4paper or letterpaper
	10pt, % Default font size, can also use 11pt or 12pt, although this is not recommended
	unnumberedsections, % Comment to enable section numbering
	twoside, % Two side traditional mode where headers and footers change between odd and even pages, comment this option to make them fixed
]{LTJournalArticle}

\addbibresource{chapitres/references.bib} % BibLaTeX bibliography file;

\runninghead{} % A shortened article title to appear in the running head, leave this command empty for no running head

\footertext{} % Text to appear in the footer, leave this command empty for no footer text

\setcounter{page}{1} % The page number of the first page, set this to a higher number if the article is to be part of an issue or larger work

%----------------------------------------------------------------------------------------
%	TITLE SECTION
%----------------------------------------------------------------------------------------

\title{\textbf{\textsc{\Huge Génération et modélisation de surfaces 3d}}} % Article title, use manual lines breaks (\\) to beautify the layout

% Authors are listed in a comma-separated list with superscript numbers indicating affiliations
% \thanks{} is used for any text that should be placed in a footnote on the first page, such as the corresponding author's email, journal acceptance dates, a copyright/license notice, keywords, etc
\author{%
	Ibrahim DOGAN, Juliette CORNU-BESSER\\ Quentin BERNARD-BOUISSIERES, Daniel DEFOING, Haluk YILMAZ \\ \textbf{Publié :} 30 avril 2026 }


% Affiliations are output in the \date{} command
\date{Projet d'année INFO-F309\\ Année Académique 2025-2025}

% Full-width abstract
\renewcommand{\maketitlehookd}{%
    \begin{abstract}
        \noindent Ce projet porte sur l’implémentation d’une application permettant de générer des cartes de bruit de Perlin en deux et en trois dimensions. Divers paramètres sont proposés afin de modifier la forme des cartes produites. 

        Le cœur du projet est développé en C++ à l’aide de la bibliothèque graphique \textbf{OpenFrameworks}. Dans le cadre des activités pédagogiques, un dispositif composé d’un bac à sable, d’un projecteur et d’un capteur Kinect est utilisé et interfacé avec l’application (\cite{ActivityApplication}). Ce système permet de reproduire physiquement une carte de bruit de Perlin bidimensionnelle en modélisant le relief dans le sable. 

        La carte de référence est ensuite comparée à celle obtenue par projection sur le bac à sable au moyen d’une application développée en Python avec la bibliothèque \textbf{Streamlit}.
    \end{abstract}
}

%----------------------------------------------------------------------------------------

\begin{document}

\maketitle % Output the title section

%----------------------------------------------------------------------------------------
%       ARTICLE CONTENTS
%----------------------------------------------------------------------------------------

\section{Introduction}

Dans le cadre du Printemps des Sciences, ce projet s’inscrit dans une démarche de génération et de modélisation de surfaces tridimensionnelles.
Il repose principalement sur l’étude du fonctionnement de l’algorithme de Perlin, ainsi que sur l’analyse de ses différents paramètres et de leurs cas limites d’application.

%------------------------------------------------

\chapter{Surfaces}

Ce chapitre se concentre sur le fonctionnement de l'algorithme de Perlin \cite{PerlinNoise}.

\section{Bruit de Perlin}

\subsection{Notion de bruit}
Lors de l'application de l'algorithme de Perlin, le bruit est considéré comme une valeur pseudo-aléatoire générée par des coordonnées données en paramètres. 
La valeur du bruit pourra être interprétée à des étapes ultérieures lors de la construction d'un espace. 
% je n'aime pas cette formulation

\subsection{Notion de fonction de bruit}

Pour mettre en relation un ensemble de bruits, les fonctions de bruit font une interpolation entre deux valeurs, soit une moyenne. 
Selon la fonction choisie, l'interpolation sera plus au moins douce.
Dans le cas de l'algorithme de Perlin, nous utilisons une interpolation linéaire entre deux points.
Le paramètre \textit{alpha} permet de jouer 
\begin{verbatim}
float lerp(float x,float y,float alpha):
	return (1 - alpha)x + alpha * y
\end{verbatim}

\subsection{Fonctionnement de l'algorithme de Perlin} 

L'algorithme de Perlin peut se décomposer en plusieurs étapes : définir l'espace, générer un échantillon fini de valeurs aléatoires, faire des produits scalaires et ensuite interpoler. 



\subsection{Application}

\subsubsection{Mouvement brownien}



\section{Extensions}

\subsection{Optimisation}

\subsection{Variantes}


%------------------------------------------------

\section{Paramètres de configuration du bruit de Perlin}

\subsection{Configuration}

\subsubsection{Amplitude}

\subsubsection{Rotation des vecteurs de gradient};

\section{Détails}

Dans cette partie, nous allons discuter sur un des paramètres qui jouent sur les détails apportés sur la carte générée.

\subsection{Fréquence}

Un des paramètres du bruit de Perlin est la fréquence dans une grille d'échantillonnage discrète.
Dans notre implémentation, nous avons appelé improprement ce paramètre \texttt{scale}.  
Ce paramètre joue sur le nombre de points dans une cellule conceptuelle.
Plus il y aura de points dans une cellule, plus les valeurs dans la cellulle donnée vont avoir des valeurs proches l'une de l'autre. 

\subsubsection{Cas limites}

Nous étudierons les différents intervalles de la fréquence, \textit{freq} pour l'évaluation de la fonction \texttt{noise (x * freq, y * freq)} où $x,y \in \mathbb{N}$. 

\subsubsection*{Fréquence prend des valeurs entières}

\begin{itemize}
	\item Les points évalués pour la fonction de bruit sont des valeurs entières, donc $\in \mathbb{Z}^2$. Or dans l'algorithme, ces coordonnées correspondent aux intersections des cellules de la grille interne.
	\item Dés lors, les vecteurs de distance calculés sont nuls et le produit scalaire est nul.
	\item L'interpolation entre valeurs nulles donne un résultat nul. Sur toute la grille, notre fonction de bruit est nulle et donne une surface monochrome puisque la fonction de bruit est constante en tout point.
	\item Cela produit une surface de couleur monochrome.
	\item Il s'agit d'un phénomène d'alignement entre la grille d'échantillonnage et la grille interne du bruit.

\end{itemize}

\begin{itemize}
		
\item Dans le cas où \texttt{freq = 0}, nous appliquons la fonction de bruit pour le même point, soit $(0,0)$.
\item Pour  le cas où \texttt{freq = 1}, nous calculons la fonction de bruit pour chaque intersection de cellules distinctes.
\item Alors que dans le cas où \texttt{freq > 1}
	\begin{itemize}
		\item Nous ne calculons pas la fonction de bruit  sur toutes les intersections mais pour toutes les $freq$ intersections de cellules.
	\end{itemize}

\end{itemize}

\subsubsection*{Fréquence prend des valeurs avec une partie décimale non nulle}

\begin{itemize}
	\item Fréquence tend vers 0, en restant strictement différent de 0
		\begin{itemize}
			\item Le domaine de la fonction de bruit est plus grand. Cela se traduit par un plus grand nombre de points pour une seule cellule donnée de la grille interne de Perlin.
			\item Pour des points voisins dans une cellule donnée, leurscoordonnées sont très proches.
			\item Les valeurs de bruit pour les différents points dans une cellules données varient très peu puisque leurs vecteurs de distance sont proches en valeurs pour les mêmes vecteurs de gradient.
			\item Dés lors, la variation de la fonction de bruit devient extrêmement faible. Ce qui donne visuellement une surface presque monochrome, proche du cas où \texttt{freq = 0}.
		\end{itemize}

	\item Fréquence prend des valeurs non entières : $freq \in \mathbb{Q}^+ \backslash \mathbb{N}$ 
	\begin{itemize}
		\item Pour les points considérés en paramètres de \texttt{noise (x * freq, y * freq)}, la fonction de bruit s'applique sur des points à l'intérieur des cellules.
		\item Dés lors, les vecteurs de distance sont différents de $\vec{0}$ et le produit scalaire est différent de $0$. 
		\item Mais si la partie fractionnaire est proche de $0$, les points considérés sont proches des coins de cellules. Les variations sont faibles et nous avons un risque de quasi-dégénérescence. 
		\item De plus, lorsque la fréquence augmente, nous utilisons moins de points et donc moins de vecteurs de gradient. Ce qui peut entrainer des problèmes de repliement spectral (\textit{aliasing}). %+ expliquer
	\end{itemize}
\end{itemize}




%------------------------------------------------

\section{Conclusion}

Ce travail a permis de présenter le fonctionnement de l’algorithme de Perlin pour la génération de cartes de bruit.
Nous avons également abordé ses extensions, notamment à travers l’utilisation des octaves et la superposition de différentes paramétrisations.
Enfin, une analyse des paramètres de configuration a été proposée, avec une attention particulière portée à la fréquence.

Dans une perspective d’approfondissement, il serait pertinent d’explorer d’autres algorithmes de génération de bruit.
%----------------------------------------------------------------------------------------
%	 REFERENCES
%----------------------------------------------------------------------------------------
\nocite{*}
\printbibliography % Output the bibliography

%----------------------------------------------------------------------------------------

\end{document}
